\documentclass[11pt]{article}
\usepackage{amsmath,amssymb,amsthm}
\usepackage[margin=1in]{geometry}
\usepackage{hyperref}

\newtheorem{theorem}{Theorem}[section]
\newtheorem{proposition}[theorem]{Proposition}
\newtheorem{lemma}[theorem]{Lemma}
\newtheorem{corollary}[theorem]{Corollary}
\theoremstyle{definition}
\newtheorem{definition}[theorem]{Definition}
\newtheorem{example}[theorem]{Example}
\theoremstyle{remark}
\newtheorem{remark}[theorem]{Remark}

\newcommand{\Stab}{\mathrm{Stab}}
\newcommand{\LO}{\mathbf{LO}}
\newcommand{\Fin}{\mathrm{Fin}}
\newcommand{\Sub}{\mathrm{Sub}}
\newcommand{\Dec}{\mathrm{Dec}}
\newcommand{\WF}{\mathrm{WF}}
\newcommand{\Acc}{\mathrm{Acc}}
\newcommand{\Hom}{\mathrm{Hom}}

\title{Dilators and ptykes as categories of embeddings,\\ constructively}
\author{Notes on a formalisation (draft, 19 November 2026)}
\date{}

\begin{document}
\maketitle

\begin{abstract}
We describe a constructive development, checked mechanically, of Girard's
dilators and ptykes~\cite{Gir81,GirPTLC} following the approach of
\emph{categories of embeddings}~\cite{Coq89}. The central result is that the category
$\Stab(C)$ of stable functors on a category of embeddings $C$ with good
finite objects is again such a category; iterating from the category of
linear orders gives dilators, ptykes of type~2, of type~3, and so on.
Well-foundedness is inductive accessibility, and no classical principle is
used. We give examples (polynomial dilators, $2^X$, $\varepsilon_X$, the
evaluation ptykes, and the Bachmann--Howard collapse $D\mapsto\vartheta D$
as a type-2 ptyx), the constructive halves of Girard's $\Pi^1_2$ and
$\Pi^1_3$ completeness theorems, and a candidate type-3 collapse aimed at
$\Pi^1_2$-comprehension.
\end{abstract}

\section{Setting}

We work in a constructive dependent type theory (with a hierarchy of
universes, function extensionality and unique choice; propositional
truncation is used for mere existence).  All
orders are \emph{decidable}: a linear order is a type $A$ with a
comparison $A\times A\to\{<,=,>\}$ satisfying reflexivity, flip,
transitivity and $x=y$ whenever the comparison says so.  A morphism of
linear orders is a strictly increasing map (an embedding).  An order is
\emph{well-founded} if every element is accessible (inductively defined
$\Acc$).  Nothing below uses excluded middle; in several places decidability
that is usually assumed is \emph{derived}.

\section{Categories of embeddings}

In~\cite{Coq89} a stable category of embeddings is the ind-completion of a
\emph{category of events}, a categorical version of Winskel's event
structures~\cite{Win}; the stable functors on it are those of Girard's
$\Pi^1_2$-logic~\cite{Gir81}. We use the following constructive variant.

\begin{definition}[Category of embeddings]
A \emph{category of embeddings} $C$ consists of a category whose morphisms
are monomorphisms, with pull-backs, together with a small type $\Fin$ of
codes and a map $\mathrm{fin}:\Fin\to C$ (the finite objects, not
necessarily injectively coded), such that any two maps from finite objects
into an object merely factor through a common finite object
(\emph{directedness}).
\end{definition}

\begin{definition}[Good finite objects]
A category of embeddings $C$ has \emph{good finite objects} (a CEF) when:
\begin{enumerate}
\item hom-sets between finite objects are Bishop-finite, with decidable
  equality;
\item each finite object has a finite list of subobjects covering, up to
  image-equivalence, every map from a finite object into it;
\item every listed subobject $s:b\to a$ has a \emph{self-amalgamation}: two
  maps $w,w':a\to d$ agreeing on $s$ whose pull-back factors through $s$;
\item finite objects are rigid (their only endomorphism is the identity);
\item finite \emph{multicoproducts} in the sense of Diers~\cite{Die80}
  (see also~\cite{Osm1,Osm2}): for finite
  objects $a,b$ a finite (junk-tolerant) list of cospans $a\to c\leftarrow
  b$ such that every pair of maps $a\to d\leftarrow b$ into a finite object
  factors through exactly one listed cospan, by exactly one map.
\end{enumerate}
\end{definition}

\begin{example}
Finite ordinals with strictly increasing maps (coded by words over
$\{0,1\}$) form a CEF: pull-backs are intersections of images, amalgamation
doubles the points outside a subobject, the multicoproducts of $[a]$ and
$[b]$ are the jointly surjective shuffles.  The category $\LO$ of all
decidable linear orders with finite objects $[n]$ is a CEF, the finiteness
data being transported from the finite ordinals.
\end{example}

\begin{definition}[Stable functors]
A \emph{stable functor} on $C$ is a functor $X:C\to\LO$ that is
\emph{continuous} (every point of $X(A)$ comes from some finite object
mapping to $A$) and \emph{stable} (it preserves pull-backs, in image
form). Morphisms of stable functors are natural transformations that are
\emph{cartesian} (naturality squares are pull-backs); they are represented
by their restriction to finite objects, so that hom-sets are small, and
extended to all objects by unique choice.
\end{definition}

In Diers's terms~\cite{Osm1}, stability says that every point factors
uniquely through a candidate; below, the candidate is the least support.

\begin{proposition}[Decidable supports]
Let $s:B\to A$ be a map with a self-amalgamation $w,w':A\to D$. For every
stable functor $X$ and $x\in X(A)$,
\[ x\in \mathrm{Im}\,X(s)\iff X(w)x=X(w')x ,\]
an equation in the linear order $X(D)$; hence membership in the image of
$X(s)$ is decidable.
\end{proposition}

This replaces the classical assumption of decidable supports (or decidable
degeneracies): it is a consequence of stability.

\begin{theorem}[Normal forms]
Let $C$ be a CEF and $X$ a stable functor. Every point $x\in X(\mathrm{fin}\,c)$
has a least support, found by a finite search over the listed subobjects,
and normal forms are unique: if $X(k)x_0=X(k')x_0$ for a point $x_0$ at its
least support then $k=k'$.
\end{theorem}

\begin{definition}[Codes]
A \emph{code} consists of $m$ generators with finite objects
$\mathrm{gen}(i)$ and, at each finite level $d$, a linear order on the
normal forms $(i,k:\mathrm{gen}(i)\to d)$, invariant under maps and
satisfying the order laws. A code $K$ \emph{decodes} to a stable functor
$\mathrm{Decode}\,K$ whose points at $A$ are pairs
$(i,k:\mathrm{fin}(\mathrm{gen}\,i)\to A)$, without any quotient.
\end{definition}

\begin{theorem}[Covering]
Every point of a stable functor $X$ is the image, under a cartesian
morphism $\mathrm{Decode}\,K\to X$, of a point of a decoded code $K$.
\end{theorem}

\begin{theorem}[Main theorem]
If $C$ is a CEF then $\Stab(C)$, the category of stable functors on $C$
and cartesian morphisms, with codes as finite objects, is again a CEF.
\end{theorem}

The proof assembles: pull-backs computed pointwise; directedness by
merging the generators of two codes, removing duplicates (decidable);
Bishop-finiteness of morphisms between decoded codes (a morphism is
determined by the images of the generators, and validity of an assignment is
decidable using the multicoproducts of $C$); subobjects as subsets of
generators; amalgamation by a lexicographic product with $\{L,R\}$;
rigidity from a constructive pigeonhole principle (an order-preserving
self-map of a Bishop-finite linear order is the identity); and the
multicoproducts of two decoded codes $K_1,K_2$: a listed cospan is given by
a subset of the generators of $K_2$ to be kept, a \emph{table} of
comparisons between generators of $K_1$ and kept generators of $K_2$ at the
listed joint types of their finite objects, which is valid (decidably), and
the resulting sum code.

\begin{corollary}[The hierarchy]
$\LO$, $\Stab(\LO)$ (pre-dilators), $\Stab(\Stab(\LO))$ (pre-ptykes of
type 2), $\Stab^3(\LO)$, \dots{} are all categories of embeddings with good
finite objects.
\end{corollary}

\section{Dilators and ptykes}

\begin{definition}
Following Girard~\cite{Gir81,GirPTLC,GN}, a pre-dilator is a stable functor
on $\LO$; it is a \emph{dilator} if
$D(A)$ is well-founded whenever $A$ is. A pre-ptyx of type~2 is a stable
functor $F$ on $\Stab(\LO)$; it is a \emph{ptyx} if $F(D)$ is
well-founded whenever $D$ is a dilator. Similarly at higher types.
\end{definition}

\begin{example}[First examples]
The identity $X\mapsto X$ is a dilator; $X\mapsto X^{\mathrm{op}}$ is a
pre-dilator that is not a dilator ($\omega^{\mathrm{op}}$ has no accessible
point). Evaluation $D\mapsto D(Z)$ at a well-founded $Z$ is a type-2 ptyx,
e.g.\ at $\omega$.
\end{example}

\begin{example}[Polynomial dilators and $2^X$]
Every polynomial built from the identity, constants, sums, products and
binary exponentiation $2^D$ is a dilator on all linear orders; in
particular $X\mapsto 2^X$ (Girard's well-ordering principle, classically
equivalent to $\mathrm{ACA}_0$~\cite{GirPTLC}), towers $2^{2^{\cdot^{\cdot^X}}}$, and the
evaluation $D\mapsto D(2^\omega)$ is a type-2 ptyx.
\end{example}

\begin{example}[$\varepsilon_X$]
Cantor normal forms over the symbols $\varepsilon_x$ ($x\in X$) form a
decidable linear order; $X\mapsto\varepsilon_X$ is a stable functor, and a
dilator: if $X$ is well-founded then so is $\varepsilon_X$. Classically
this principle is equivalent to $\mathrm{ACA}_0^+$ (Afshari--Rathjen~\cite{AR},
Marcone--Montalb\'an~\cite{MM}; see also the survey~\cite{Rat21}).
\end{example}

\begin{theorem}[Bachmann--Howard collapse]
For every pre-dilator $D$, Freund's order $\vartheta D$~\cite{Fre19} (terms
$\vartheta(\sigma)$ with $\sigma\in D(\vartheta D)$, compared by
\[\vartheta\sigma<\vartheta\tau\iff \vartheta\sigma\le\text{some support
element of }\tau,\ \text{or}\ \sigma<\tau\text{ and all support elements of
}\sigma\text{ are }<\vartheta\tau)\]
is a decidable linear order; $D\mapsto\vartheta D$ is a stable functor on
$\Stab(\LO)$; and it is a type-2 ptyx: $\vartheta D$ is well-founded for
every dilator $D$.
\end{theorem}

Classically, ``$\vartheta D$ is well-founded for every dilator $D$'' is
equivalent to $\Pi^1_1$-comprehension (Freund~\cite{Fre19}; see
also~\cite{AFW} for Girard's $\Lambda$ and P\"appinghaus's functor~$J$). The constructive proof
takes the accessible part $W$ of $\vartheta D$, uses that $D(W)$ is
well-founded, and argues by induction on the image of $\sigma$ in $D(W)$,
with a semantic lemma identifying the term comparison with the comparison
in $D$ of the image of a well-behaved sub-order.

\section{Girard's completeness theorems, constructive halves}

Both theorems are due to Girard~\cite{Gir81,GirPTLC}; we formalise the
directions that need no classical principle.

\begin{theorem}[$\Pi^1_2$]
To a tree $T$ (a $\Pi^1_2$ statement $\forall f\exists g\forall n\,R$)
one associates a pre-dilator $D_T$ (labelled trees ordered by the
Kleene--Brouwer order) such that $D_T$ is a dilator iff all its attempt
trees are well-founded, and a rank of some $U_f$ into a well-founded order
refutes that $D_T$ is a dilator.
\end{theorem}

\begin{theorem}[$\Pi^1_3$]
The same construction one type up gives, from a $\Pi^1_3$ statement, a
type-2 pre-ptyx (trees whose nodes carry partial natural embeddings into
the dilator argument, compared on the listed joint types), which is a ptyx
iff the corresponding trees are well-founded for every dilator.
\end{theorem}

\begin{theorem}[$\Pi^1_4$]
Fix a sequence of codes of finite dilators, each decoding to a dilator. A
$\Pi^1_4$ statement $\forall f\,\exists g\,\forall h\,\exists k\ldots$ gives a
type-3 pre-ptyx $G$: its value at a type-2 pre-ptyx $F$ is the
Kleene--Brouwer order of finite attempts to embed, naturally, the finite
part of the type-2 pre-ptyx $F_u$ of the previous theorem into $F$
(comparisons checked on the listed multicoproducts of codes). Then $G$ is a
type-3 ptyx iff every 2-ptyx has a well-founded attempt tree; and if $F_f$
is a 2-ptyx for some $f$, then $G$ is not a type-3 ptyx.
\end{theorem}

The converse directions use bar induction and are not formalised. A
positive instance (the tree of strictly decreasing sequences) is proved to
give a dilator.

The hypothesis that the listed codes decode to dilators is discharged
for \emph{monotone} codes (Girard's monotonicity of predilators on the
generators: $k\le k'$ pointwise implies $(i,k)\le(i,k')$). Not every code
is a dilator: $X\mapsto X^{\mathrm{op}}$ is a one-generator code which is
not.

\begin{theorem}[Monotone finite predilators are dilators]
If a code $K$ is monotone then $\mathrm{Decode}\,K$ is a dilator.
\end{theorem}

The proof uses almost-full relations: following Coquand's direct proof of
the intuitionistic Ramsey theorem~\cite{CoqRamsey}, if a tree $p$ secures
$A\vee R$ and $q$ secures $B\vee S$ for $n$-ary $R,S$, then $p\otimes_n q$
secures $A\vee B\vee(R\wedge S)$; hence almost-full relations are closed
under intersection. On a well-founded linear order, $\le$ is almost full;
almost-fullness pulls back along maps; so ``same generator and pointwise
$\le$'' is almost full on $\mathrm{Decode}\,K(A)$; by monotonicity it is
contained in $\le$, and an almost-full relation contained in $\le$ makes a
linear order well-founded.

\begin{corollary}[A statement beyond the type theory]
Take for $S$ the $\Pi^1_4$ axiom of $\Pi^1_2$-$\mathrm{CA}_0$ and let $G_S$ be
the type-3 pre-ptyx of the previous theorem. The formalised implication
``$G_S$ is a type-3 ptyx $\Rightarrow$ no $F_f$ is a 2-ptyx'', followed by
the classical type-2 completeness theorem and Kleene's normal form, gives
\[ G_S\ \text{is a type-3 ptyx}\ \Longrightarrow\ \Pi^1_2\text{-}\mathrm{CA}_0 .\]
The statement ``$G_S$ is a type-3 ptyx'' is expressible in the type theory
used here. Since this type theory is predicative, of strength well below
$\Pi^1_2$-$\mathrm{CA}_0$, the statement is not provable in it (a
metatheorem, via G\"odel's second incompleteness theorem; not formalised).
\end{corollary}

\section{What one may hope for}

\paragraph{A type-3 collapse.} Girard~\cite{GirH} conjectured that
$\Pi^1_n$-comprehension is equivalent to the totality of a canonical
``exponential'' $A_n$ of type $n+1$; known for $n=0$ ($2^X$) and $n=1$
($\Lambda$, or Freund's $\vartheta$~\cite{Fre19,AFW}). Freund showed that
fixed points of normal dilators correspond to $\Pi^1_1$-induction, and fixed
points of normal 2-ptykes to $\Pi^1_2$-induction~\cite{FrePilot}, while
Bachmann--Howard \emph{collapses} of dilators correspond to
$\Pi^1_1$-comprehension~\cite{Fre19}. The
natural candidate for $n=2$ is therefore the \emph{collapse of a 2-ptyx}:
for $P:\mathrm{DIL}\to\mathrm{DIL}$, a dilator $\Theta_P$ with a collapse
$\vartheta:P(\Theta_P)\Rightarrow\Theta_P$, so that $P$ is applied to the
dilator being built. Its terms are finite: $\vartheta(K,z,\vec y,\vec t)$
with $K$ a code, $z\in P(\mathrm{Decode}\,K)([c])$, $\vec y$ the arguments
and $\vec t$ the images of the generators of $K$ (its support). The
supports are \emph{functions} (trace terms with variables), and, following
Towsner~\cite{TowPoly}, two collapses are compared by a closure rule:
$\vartheta\sigma<\vartheta\tau$ if $\vartheta\sigma$ is not closed under
some support function $\gamma$ of $\tau$, i.e.\ $\gamma$ with its variables
placed just below $\vartheta\sigma$ is $\ge\vartheta\sigma$; symmetrically;
and otherwise $\sigma$ and $\tau$ are compared inside $P$ after selecting
the amalgam of the two codes from the multicoproduct list. Atoms lie below
all collapses; placing them above (as cardinals) is incompatible with the
closure rule for a total order. The
statement ``$\Theta_P$ is a dilator for every 2-ptyx $P$'' is $\Pi^1_4$,
the complexity of $\Pi^1_2$-$\mathrm{CA}_0$; no well-ordering principle for
$\Pi^1_2$-$\mathrm{CA}_0$ is known (compare Rathjen's conjecture on
$\beta$-models~\cite[Conj.~6.6]{Rat21}, a type-2 statement, and the ordinal
analysis~\cite{Rat05}).

What is established: a naive variant (collapses without free variables)
reduces to $\vartheta$ of an ordinary dilator and is provable in
$\Pi^1_1$-$\mathrm{CA}_0$; the candidate reduces to $\vartheta_D$ when $P$
lifts $D$; a single instance (iteration of $\Delta$ along $\omega$) should
exceed $\Pi^1_1$-$\mathrm{CA}_0$, modulo a comparison of Freund's
$\vartheta$-order with Buchholz's $\psi$-order~\cite{Buc86}.  Open: an informal proof
that $\Theta_P$ is well-founded for every 2-ptyx (some form of induction on
dilators), and the converse.

\paragraph{Towsner's polymorphic notations.} Towsner's ordinal
analyses~\cite{TowPMP,TowHalf,TowPoly} towards parameter-free
$\Pi^1_2$-comprehension use ordinal bounds with free
variables (functions that preserve well-foundedness, i.e.\ dilators) and
cardinals of function sort, whose coefficient sets contain functions that
are instantiated at the parameters of the other term. These are dilators
and type-2 ptykes in all but name; the semantics above could give them a
precise constructive meaning. His rule for function-valued coefficients
gives the closure rule used above; his placement of collapses below
variables, on the other hand, relies on his order being partial on open
terms and is not available for a dilator, whose order is total.

\paragraph{Other well-ordering principles.} $X\mapsto\varphi X0$
($\mathrm{ATR}_0$) and $X\mapsto\Gamma_X$ fit the same pattern as
$\varepsilon_X$.

\section{What has been implemented}

The development is written in a small dependently typed language and
checked by its type checker (about 26{,}800 declarations for the whole
library; the files below take a few seconds each). Main files, in the
\texttt{UNIT/} directory:
\begin{itemize}
\item categories of embeddings: \texttt{CatEmb}, \texttt{CatEmbFin} (CEF,
  Diers multicoproducts), \texttt{StabFun}, \texttt{StabPb},
  \texttt{StabDec} (decidable supports), \texttt{StabSupp} (least supports,
  unique normal forms), \texttt{StabCode}, \texttt{StabCover},
  \texttt{StabDir}, \texttt{StabFin}, \texttt{StabHom}, \texttt{StabSub},
  \texttt{StabAmal}, \texttt{StabRigid}, \texttt{StabJt0},
  \texttt{StabSum}, \texttt{StabJt}, \texttt{StabCEF} (main theorem);
\item base instances: \texttt{FinOrd}, \texttt{FinOrdEx},
  \texttt{LOCat}, \texttt{LOFin}, \texttt{LOCEF};
  well-foundedness hierarchy: \texttt{LOWF};
\item Girard's reductions: \texttt{KBTree}, \texttt{LabTree},
  \texttt{Girard}, \texttt{GirardPos}, \texttt{LabTree2}, \texttt{Girard2},
  and at type 3 \texttt{NSumLO}, \texttt{LabTree3}, \texttt{Girard3};
\item almost-full relations and the intuitionistic Ramsey theorem:
  \texttt{AlmostFull}; monotone finite predilators are dilators:
  \texttt{FinDilWF};
\item examples: \texttt{DilLO}, \texttt{LOEx} (polynomials, $2^X$,
  towers, evaluations), \texttt{EpsTerm}, \texttt{EpsFun}, \texttt{EpsWF}
  ($\varepsilon_X$), \texttt{ThetaTerm}, \texttt{ThetaLO},
  \texttt{ThetaLO2}, \texttt{ThetaLO3}, \texttt{ThetaFun},
  \texttt{ThetaWF}, \texttt{ThetaFun2} ($\vartheta$ as a type-2 ptyx);
\item type 3 (in progress): \texttt{ThetaP3Term} (terms of $\Theta_P$ and a
  decidable comparison, defined with a fuel bound); a Python mirror of the candidate in
  \texttt{mirror/}.
\end{itemize}
An earlier, more hands-on development over finite ordinals
(\texttt{CatDil}, \texttt{DilPtyx}, \texttt{Ptyx*}, \texttt{Ptyx3*}) treats
ptykes of types 2 and 3 directly, with germs, subptykes and directed
colimits; it is subsumed by the abstract theorem.

\begin{thebibliography}{99}
\bibitem{AR} B.~Afshari and M.~Rathjen, Reverse mathematics and
  well-ordering principles: a pilot study, \emph{Ann.\ Pure Appl.\ Logic}
  160 (2009) 231--237.
\bibitem{AFW} J.~P.~Aguilera, A.~Freund and A.~Weiermann, Induction on
  dilators and Bachmann--Howard fixed points, arXiv:2412.13051.
\bibitem{Buc86} W.~Buchholz, A new system of proof-theoretic ordinal
  functions, \emph{Ann.\ Pure Appl.\ Logic} 32 (1986) 195--207.
\bibitem{CoqRamsey} Th.~Coquand, A direct proof of Ramsey's Theorem, note,
  2010, \url{https://www.cse.chalmers.se/~coquand/ramsey1.pdf}.
\bibitem{Coq89} Th.~Coquand, Categories of embeddings, \emph{Theoret.\
  Comput.\ Sci.} 68 (1989) 221--237.
\bibitem{Die80} Y.~Diers, Cat\'egories localement multipr\'esentables,
  \emph{Arch.\ Math.} 34 (1980) 344--356.
\bibitem{Fre19} A.~Freund, $\Pi^1_1$-comprehension as a well-ordering
  principle, \emph{Adv.\ Math.} 355 (2019) 106767.
\bibitem{FrePilot} A.~Freund, Well ordering principles and
  $\Pi^1_4$-statements: a pilot study, arXiv:2006.12111.
\bibitem{Gir81} J.-Y.~Girard, $\Pi^1_2$-logic, Part 1: Dilators,
  \emph{Ann.\ Math.\ Logic} 21 (1981) 75--219.
\bibitem{GirH} J.-Y.~Girard, Herbrand's theorem and proof theory, in:
  \emph{Proceedings of the Herbrand Symposium (Marseille 1981)},
  North-Holland, 1982, 29--38.
\bibitem{GirPTLC} J.-Y.~Girard, \emph{Proof Theory and Logical
  Complexity}, vol.~1, Bibliopolis, 1987; vol.~2, manuscript.
\bibitem{GN} J.-Y.~Girard and D.~Normann, Embeddability of ptykes,
  \emph{J.\ Symbolic Logic} 57 (1992) 659--676.
\bibitem{MM} A.~Marcone and A.~Montalb\'an, The Veblen functions for
  computability theorists, \emph{J.\ Symbolic Logic} 76 (2011) 575--602.
\bibitem{Osm1} A.~Osmond, On Diers's theory of spectrum I: stable functors
  and right multi-adjoints, arXiv:2012.00853.
\bibitem{Osm2} A.~Osmond, On Diers's theory of spectrum II: geometries and
  dualities, arXiv:2012.02167.
\bibitem{Rat05} M.~Rathjen, An ordinal analysis of parameter free
  $\Pi^1_2$-comprehension, \emph{Arch.\ Math.\ Logic} 44 (2005) 263--362.
\bibitem{Rat21} M.~Rathjen, Well-ordering principles in proof theory and
  reverse mathematics, arXiv:2010.12453.
\bibitem{TowPMP} H.~Towsner, Proofs that modify proofs, arXiv:2403.17922.
\bibitem{TowHalf} H.~Towsner, Proofs that modify proofs, 1/2,
  arXiv:2506.16491.
\bibitem{TowPoly} H.~Towsner, Polymorphic ordinal notations,
  arXiv:2504.02131.
\bibitem{Win} G.~Winskel, Event structures, Cambridge University Computer
  Laboratory, Technical Report No.~95.
\end{thebibliography}

\end{document}
